\documentclass[10pt,a4paper]{amsart}
\usepackage[margin=19mm]{geometry}
\usepackage{amsmath,amssymb,mathtools}
\usepackage[scr=rsfs,cal=euler]{mathalfa}
\usepackage{xfrac}
\usepackage[unicode,hidelinks,bookmarks=false]{hyperref}
\newcommand{\ie}{i.e.\ }
\newcommand{\cf}{cf.\ }
\newcommand{\resp}{respectively\ }
\newcommand{\Subsection}{\subsection}
\providecommand{\nobreakdash}{\nobreak\hbox{-}}
\setlength{\emergencystretch}{2em}
\multlinegap0pt
\usepackage{amssymb}
\usepackage{xcolor}
\usepackage{algorithm}\usepackage{algpseudocode}
\newtheorem{theorem}{Theorem}
\title{A Note on Quaternions over Commutative Rings}
\author{Antony Telveenus}
\date{Source: arXiv:2606.23376v3 (CC BY 4.0)}
\begin{document}
\maketitle

\noindent\emph{Source and licence.} A Note on Quaternions over Commutative Rings. Version arXiv:2606.23376v3 (CC BY 4.0), distributed by arXiv under the Creative Commons Attribution 4.0 International licence, http://creativecommons.org/licenses/by/4.0/, as recorded in the arXiv licence metadata for this exact version and shown on the abstract page https://arxiv.org/abs/2606.23376v3. Full licence notice: \texttt{licenses/arxiv-2606.23376-CC-BY-4.0.txt}.\par\smallskip

\noindent\textit{Editorial scope.} Counted unit: the article's theorem theorem (source0.tex lines 78--80). The article's complete original proof of it is delivered with it (source0.tex lines 81--133, one \texttt{proof} environment of 53 lines). No mathematics is written, repaired or re-derived: the statement and the proof are the article's own lines, in the article's own notation, and no retained line is substituted or re-flowed.  No further source-local context is needed: every cross-reference of every retained line resolves inside the two delivered spans.  Nothing was omitted from any retained span, so the package carries no omission record.  No retained line contains a citation, so the wrapper carries no bibliography and no bibliography entry.  No retained line contains a figure environment, an image-inclusion command or drawing code, so no image is delivered and the package declares no asset dependency.  Editorial formatting only, in addition: an AMS \texttt{amsart} preamble that restates the article's own declarations and macros used by the retained lines, namely the environments \texttt{theorem} on separate counters, together with the standard packages those lines need.  The retained environments print in this document as Theorem 1 in order of appearance, whereas the article prints the same items with the numbering of its own full document.  The complete native source of the article as distributed in the arXiv e-print for this exact version is bundled unchanged as \texttt{sources/203425/source0.tex}.
\begin{theorem}
Let $\mathcal{H_{\mathbf{R}}}=\{ q=w+xi+yj+zk \ | w,x,y,z\in \mathbf{R}\}$ be the ring of quaternions over $\mathbf{R}$. Then there exists an induced ring of quaternions $\mathcal{H_{\mathbf{R}}^{\perp}}$ with orientation just opposite to $\mathcal{H_{\mathbf{R}}}$ such that $\mathcal{H}_{\mathbf{R}}\cong\mathcal{H_{\mathbf{R}}^{\perp}}$. 
\end{theorem}

\begin{proof}
We define $\mathcal{H_{\mathbf{R}}^{\perp}}=\{ Q=wI+x\Gamma_{1}+y\Gamma_{2}+z\Gamma_{3} \ | w,x,y,z\in \mathbf{R}\}$, where \\
$I=\left(
\begin{array}{cccc}
1 & 0 & 0 & 0 \\
0 & 1 & 0 & 0 \\
0 & 0 & 1 & 0 \\
0 & 0 & 0 & 1 \\
\end{array}
\right)
$,
$\Gamma_{1}=\left(
\begin{array}{cccc}
0 & -1 & 0 & 0 \\
1 & 0 & 0 & 0 \\
0 & 0 & 0 & 1 \\
0 & 0 & -1 & 0 \\
\end{array}
\right)
$,
$\Gamma_{2}=\left(
\begin{array}{cccc}
0 & 0 & -1 & 0 \\
0 & 0 & 0 & -1 \\
1 & 0 & 0 & 0 \\
0 & 1 & 0 & 0 \\
\end{array}
\right)
$,
$\Gamma_{3}=\left(
\begin{array}{cccc}
0 & 0 & 0 & -1 \\
0 & 0 & 1 & 0 \\
0 & -1 & 0 & 0 \\
1 & 0 & 0 & 0 \\
\end{array}
\right)
$.
A bit lengthy calculation shows that $\mathcal{H_{\mathbf{R}}^{\perp}}$ satisfies all rules of addition and multiplication except the interchange of rule 2 and rule 3 of multiplication. 
Let $f:\mathcal{H_{\mathbf{R}}} \mapsto \mathcal{H_{\mathbf{R}}^{\perp}}$ be defined by $f(w+xi+yj+zk)=wI-x\Gamma_{1}-y\Gamma_{2}-z\Gamma_{3}$. This map is clearly a bijective map preserving the operation  addition and $f(1)=I$. Let us show that f preserves the multiplication as well. \\
Let $q_{1}=w_{1}+x_{1}i+y_{1}j+z_{1}k$, $q_{2}=w_{2}+x_{2}i+y_{2}j+z_{2}k$ be any two elements in $\mathcal{H}$. Then $f(q_{1})=w_{1}I-x_{1}\Gamma_{1}-y_{1}\Gamma_{2}-z_{1}\Gamma_{3}$ and $f(q_{2})=w_{2}I-x_{2}\Gamma_{1}-y_{2}\Gamma_{2}-z_{2}\Gamma_{3}$. \\
$q_{1}q_{2}=(w_{1}w_{2}-x_{}x_{2}-y_{1}y_{2}-z_{1}z_{2})+(w_{1}x_{2}+w_{2}x_{1}+y_{1}z_{2}-z_{1}y_{2})i+(w_{1}y_{2}+w_{2}y_{1}-x_{1}z_{2}+z_{1}x_{2})j
+(w_{1}z_{2}+z_{1}w_{2}+x_{1}y_{2}-x_{2}y_{1})k.$ \\
$f(q_{1}q_{2})=(w_{1}w_{2}-x_{}x_{2}-y_{1}y_{2}-z_{1}z_{2})I-(w_{1}x_{2}+w_{2}x_{1}+y_{1}z_{2}-z_{1}y_{2})\Gamma_{1}-(w_{1}y_{2}+w_{2}y_{1}-x_{1}z_{2}+z_{1}x_{2})\Gamma_{2}
+(w_{1}z_{2}+z_{1}w_{2}+x_{1}y_{2}-x_{2}y_{1})\Gamma_{3}.$ \\
\noindent
Using the relations $\Gamma_{1}^{2}=\Gamma_{2}^{2}=\Gamma_{3}^{2}=-1,\Gamma_{1}\Gamma_{2}=-\Gamma_{3},\Gamma_{2}\Gamma_{3}=-\Gamma_{1},\Gamma_{3}\Gamma_{1}=-\Gamma_{2},
\Gamma_{2}\Gamma_{1}=\Gamma_{3},\Gamma_{3}\Gamma_{1}=\Gamma_{1},\Gamma_{1}\Gamma_{3}=\Gamma_{2}$, we have 
$f(q_{1})f(q_{2})=(w_{1}I-x_{1}\Gamma_{1}-y_{1}\Gamma_{2}-z_{1}\Gamma_{3})(w_{2}I-x_{2}\Gamma_{1}-y_{2}\Gamma_{2}-z_{2}\Gamma_{3})=(w_{1}w_{2}-x_{1}x_{2}-y_{1}y_{2}-z_{1}z_{2})I
-(w_{1}x_{2}+w_{2}x_{1}+y_{1}z_{2}-z_{1}y_{2})\Gamma_{1}-(w_{1}y_{2}+w_{2}y_{1}-x_{1}z_{2}+z_{1}x_{2})\Gamma_{2}
+(w_{1}z_{2}+z_{1}w_{2}+x_{1}y_{2}-x_{2}y_{1})\Gamma_{3}=f(q_{1}q_{2}).$  \\
$ \therefore \quad \quad \mathcal{H_{\mathbf{R}}}\cong\mathcal{H_{\mathbf{R}}^{\perp}} $
\end{proof}

\end{document}
